\documentclass[11pt]{article}
\usepackage{series}

\hypersetup{
  unicode=true,
  pdftitle={Entanglement, Locality, and Measurement from Coherent Sector Dynamics},
  pdfauthor={Peter Nero}
}

\title{Entanglement, Locality, and Measurement from Coherent Sector Dynamics}
\author{Peter Nero}
\date{Janury 2026}

\begin{document}
\maketitle


\begin{abstract}
Algebraic quantum field theory (AQFT) enforces relativistic locality at the level of observable
algebras while allowing globally entangled states. However, AQFT alone does not explain why
physically realized states are typically entangled, how entanglement propagates, or why
measurement reduces entanglement without enabling superluminal signaling. We present a local,
state-restricting framework in which physically admissible states are those induced by coherent-sector
dynamics (bounded geometry, spectral gap, bounded coherent projector, and contractive stability).
Entanglement is shown to arise as a global coherence constraint in configuration space, while
microcausality and the time-slice property remain intact on spacetime. Measurement is modeled as
localized disturbance followed by stabilization into admissible basins, yielding partial
disentanglement. Standard Bell/CHSH and temporal Bell (Leggett--Garg) violations are recovered
without nonlocal influence. Entanglement spreading via common ancestors (mediated entanglement and
entanglement swapping) is interpreted as causal propagation of coherence constraints through local
interactions. The framework integrates the MTT QFT-projection layer, the measurement-as-stabilization
model, and the Bell/temporal-Bell analyses in a single AQFT-compatible description.
\end{abstract}

\tableofcontents

% ============================================================
\section{Introduction}

Quantum entanglement is often presented as a tension between locality and correlation.
Relativistic quantum field theory resolves this tension formally in the algebraic framework
(AQFT), which enforces strict locality at the level of observable algebras while allowing
highly entangled global states \cite{Haag1996,BFV2003,HollandsWald2001}.
AQFT thereby explains why entanglement does not permit superluminal signaling, but it does not
explain why physically realized states are typically entangled, how entanglement spreads, or how
measurement reduces entanglement.

In Modal Triplet Theory (MTT), physical states are not arbitrary states on the AQFT net:
they are restricted by coherence admissibility (bounded geometry, spectral gap, bounded coherent
projector, and stability margins encoded by the Fundamental Contractivity Condition, FCC)
\cite{MTTFoundation,FP3,FP5,FP6}.
Measurement is modeled as localized disturbance followed by stabilization into an admissible
basin \cite{NeroMeasurement}, and Bell/temporal-Bell violations are treated as projection
artifacts of globally consistent coherent histories \cite{NeroBellBeables,NeroTemporalBell}.

This paper provides a unified, AQFT-compatible formulation:
\begin{itemize}
\item We do \emph{not} modify AQFT axioms (locality, isotony, time-slice).
\item We restrict the physically relevant state space to an admissible class
      $\mathcal{S}_{\mathrm{coh}}\subset\mathcal{S}(\mathcal{A})$ induced by coherent-sector dynamics.
\item Entanglement is explained as a coherence constraint in configuration space, while microcausality
      and causal propagation remain intact on spacetime.
\item Measurement reduces entanglement by disturbance + stabilization, without nonlocal influence.
\item Entanglement propagation through common ancestors is treated as standard local interaction and
      causal spreading (Lieb--Robinson-type bounds).
\end{itemize}

% ============================================================
\section{AQFT preliminaries: locality and states}

Let $(Y^4,g)$ be a globally hyperbolic spacetime. AQFT assigns to each suitable region
$\mathcal{O}\subset Y^4$ a C$^\ast$-algebra $\mathcal{A}(\mathcal{O})$ such that:

\paragraph{Isotony.}
If $\mathcal{O}_1\subset\mathcal{O}_2$, then $\mathcal{A}(\mathcal{O}_1)\subset\mathcal{A}(\mathcal{O}_2)$.

\paragraph{Microcausality (local commutativity).}
If $\mathcal{O}_1$ and $\mathcal{O}_2$ are spacelike separated, then
\begin{equation}\label{eq:microcausality}
[\mathcal{A}(\mathcal{O}_1),\mathcal{A}(\mathcal{O}_2)]=0.
\end{equation}

\paragraph{Time-slice axiom.}
If $\mathcal{O}$ contains a Cauchy surface for its domain of dependence, then
$\mathcal{A}(\mathcal{O})$ generates the algebra in the domain of dependence.

A \emph{state} is a positive normalized linear functional
\[
\omega:\mathcal{A}\to\mathbb{C},\qquad \omega(\mathbf{1})=1,\quad \omega(A^\ast A)\ge 0.
\]
States need not factorize over commuting subalgebras, and AQFT explicitly permits entangled states
\cite{Haag1996,BFV2003}.

% ============================================================
\section{Entanglement in AQFT and the factorization gap}

Let $\mathcal{A}_1:=\mathcal{A}(\mathcal{O}_1)$ and $\mathcal{A}_2:=\mathcal{A}(\mathcal{O}_2)$ for spacelike-separated regions.
Microcausality implies $[A,B]=0$ for $A\in\mathcal{A}_1$ and $B\in\mathcal{A}_2$, but it does \emph{not} imply state factorization.

A state is \emph{separable} on $\mathcal{A}_1\vee\mathcal{A}_2$ if it admits a convex decomposition
\begin{equation}\label{eq:separable}
\omega(AB)=\sum_i p_i\,\omega^{(1)}_i(A)\,\omega^{(2)}_i(B),
\quad A\in\mathcal{A}_1,\;B\in\mathcal{A}_2,\; p_i\ge 0,\;\sum_i p_i=1.
\end{equation}
Otherwise it is entangled. Thus entanglement is compatible with AQFT locality:
non-factorization is a property of the state, not of commutators.

AQFT therefore answers the \emph{compatibility} question (“how can entanglement exist without signaling?”),
but it does not by itself supply a selection principle for why physically realized states are typically non-factorizing.

% ============================================================
\section{Coherent-sector admissibility as a state-selection principle}

\subsection{Two-layer picture: configuration versus observables}

MTT distinguishes:
\begin{enumerate}
\item a higher-dimensional coherent configuration layer (10D / bundle manifold) in which admissibility and stability
      are formulated (spectral gaps, projector boundedness, contractive evolution);
\item an effective 4D observable layer described by AQFT on $(Y^4,g)$.
\end{enumerate}
The QFT-projection paper formalizes the effective AQFT output of coherent-sector reduction and renormalization on curved
spacetimes (Hadamard states, local covariance, and standard renormalization freedoms) \cite{NeroMTTtoQFT,Haag1996,BFV2003,HollandsWald2001}.

\subsection{Admissible state class}

\begin{definition}[Admissible (coherent) state class]\label{def:Scoh}
Let $\mathcal{S}(\mathcal{A})$ denote the space of AQFT states. Define $\mathcal{S}_{\mathrm{coh}}\subset\mathcal{S}(\mathcal{A})$
to be the subclass of states induced by admissible coherent configurations, i.e. those satisfying:
\begin{enumerate}[label=\textup{(\roman*)}]
\item bounded geometry on time slabs;
\item persistence of a uniform spectral gap separating coherent from noncoherent modes;
\item boundedness of the coherent projector on Sobolev scales;
\item local disturbance--damping stability margins (FCC) on the relevant invariant sets.
\end{enumerate}
\end{definition}

This is a state-selection principle: AQFT axioms remain unchanged, but the set of physically realized states is restricted.

\begin{proposition}[Locality preserved]\label{prop:locality-preserved}
For any $\omega\in\mathcal{S}_{\mathrm{coh}}$, microcausality \eqref{eq:microcausality} holds exactly as in AQFT.
\end{proposition}

\begin{proof}
Microcausality is an algebraic property of the net $\mathcal{O}\mapsto\mathcal{A}(\mathcal{O})$, independent of state choice.
Restricting to $\mathcal{S}_{\mathrm{coh}}$ does not alter commutators.
\end{proof}

\subsection{Entanglement as coherence constraint}

\begin{proposition}[Entanglement from global coherence constraints]\label{prop:entanglement-coherence}
States in $\mathcal{S}_{\mathrm{coh}}$ may be non-factorizing over spacelike-separated regions. Such non-factorization is interpreted
as a coherence constraint in configuration space, not as superluminal influence in spacetime.
\end{proposition}

\begin{remark}
The phrase “global” here refers to configuration-space admissibility and stability constraints, not to a physical
signal propagating in spacetime. This is the same logical distinction used throughout FP~V--VI to separate global
attractors in configuration space from unique spacetime histories \cite{FP5,FP6}.
\end{remark}

% ============================================================
\section{Bell/CHSH violations without nonlocal influence}

Let $A_0,A_1\in\mathcal{A}(\mathcal{O}_A)$ and $B_0,B_1\in\mathcal{A}(\mathcal{O}_B)$ be observables with spectra in $[-1,1]$,
with $\mathcal{O}_A$ spacelike to $\mathcal{O}_B$.
Define the CHSH expression
\begin{equation}\label{eq:CHSH}
S:=\omega(A_0B_0)+\omega(A_0B_1)+\omega(A_1B_0)-\omega(A_1B_1).
\end{equation}

Classical local hidden-variable models yield $|S|\le 2$ \cite{Bell1964,CHSH1969}. Quantum theory permits $|S|\le 2\sqrt{2}$
(Tsirelson bound) \cite{Tsirelson1980}.

\begin{proposition}[Bell violation within AQFT locality]\label{prop:bell}
For $\omega\in\mathcal{S}_{\mathrm{coh}}$, CHSH inequalities can be violated while microcausality holds.
\end{proposition}

\begin{proof}[Explanation]
Microcausality ensures $[A_i,B_j]=0$ but does not imply separability \eqref{eq:separable}.
CHSH derivations of $|S|\le 2$ require a factorizable hidden-variable representation, which fails for generic
entangled states. Since $\mathcal{S}_{\mathrm{coh}}$ includes non-factorizing admissible states, violations are possible.
No signaling follows because commutators vanish and local marginals cannot be controlled by spacelike choices.
\end{proof}

This connects directly to the MTT Bell analysis: Bell nonlocality is a failure of classical factorization,
interpreted as a projection artifact of a local higher-dimensional ontology \cite{NeroBellBeables}.

% ============================================================
\section{Temporal Bell (Leggett--Garg) violations from disturbance + stabilization}

Let $Q(t)$ be a dichotomic observable ($\pm 1$) associated with a system measured at times $t_1<t_2<t_3$.
Define correlators $C_{ij}:=\omega(Q(t_i)Q(t_j))$. A standard Leggett--Garg inequality is \cite{LeggettGarg1985}
\begin{equation}\label{eq:LG}
K := C_{12}+C_{23}-C_{13}\le 1
\end{equation}
under macrorealism and noninvasive measurability.

\begin{proposition}[Temporal Bell violations]\label{prop:temporal}
Temporal Bell/Leggett--Garg violations in $\mathcal{S}_{\mathrm{coh}}$ arise because measurement is necessarily invasive:
it acts as a disturbance that perturbs coherence, followed by stabilization into an admissible basin.
\end{proposition}

\begin{remark}
This is precisely the mechanism formalized in the temporal Bell and measurement papers:
violations reflect global consistency constraints on admissible histories together with unavoidable disturbance \cite{NeroTemporalBell,NeroMeasurement}.
\end{remark}

% ============================================================
\section{Measurement as disturbance and stabilization: partial disentanglement}

We model measurement as a localized interaction channel described by a completely positive (CP) instrument.
In density-operator language, an outcome channel $m$ is
\begin{equation}\label{eq:instrument}
\rho \mapsto \rho_m'=\frac{\mathcal{M}_m(\rho)}{\Tr(\mathcal{M}_m(\rho))},
\qquad
\mathcal{M}_m(\rho)=\sum_\alpha K_{m,\alpha}\rho K_{m,\alpha}^\ast.
\end{equation}
In MTT interpretation, the Kraus operators represent:
(i) a localized disturbance (kicking modal amplitudes), and
(ii) stabilization/damping returning the configuration to an admissible basin \cite{NeroMeasurement}.

\begin{proposition}[Partial disentanglement]\label{prop:partial}
Let $\rho_{AB}$ be a bipartite state. A local measurement on $A$ produces
\[
\rho_{AB}\mapsto \rho'_{AB,m}=\frac{(\mathcal{M}_m\otimes I)\rho_{AB}}{\Tr((\mathcal{M}_m\otimes I)\rho_{AB})},
\]
which generically reduces entanglement between $A$ and $B$ by damping off-diagonal coherence
in the measured sector.
\end{proposition}

\begin{remark}
This explains “collapse-like” behavior without superluminal influence: the operation is local,
but the global admissible state is updated after disturbance and stabilization.
\end{remark}

% ============================================================
\section{Entanglement propagation: common ancestors and causal spreading}

Entanglement can be created between subsystems that never directly interacted by local interactions with a third system.

\subsection{Mediated entanglement and entanglement swapping}

Let $A,B,C$ be three subsystems. Suppose $A$ interacts locally with $C$ and later $B$ interacts locally with $C$.
After tracing out $C$, $\rho_{AB}$ can become entangled even if $A$ and $B$ had no shared creation history.

In MTT terms, $C$ acts as a \emph{common coherence ancestor} in configuration space: local interactions impose a joint
admissibility constraint that propagates through admissible dynamics.

\subsection{Finite-speed spreading}

In many-body and lattice models, locality of interactions implies Lieb--Robinson bounds \cite{LiebRobinson1972}:
for observables $A$ and $B$ supported at distance $d$,
\begin{equation}\label{eq:LR}
\|[A(t),B]\|\le C\,e^{-\mu(d-vt)}.
\end{equation}
This expresses finite-speed propagation of influence/correlation. In the MTT coherent regime,
effective causal cones and bounded-geometry dynamics provide analogous control: correlations spread through local interaction,
but signaling remains bounded by the effective light cone.

\begin{remark}
Equation \eqref{eq:LR} is a standard template; the precise constants depend on the effective dynamics.
The conceptual point is that “entanglement spreading” is causal: it proceeds through a chain of local interactions
and does not constitute superluminal influence.
\end{remark}

% ============================================================
\section{What is explained and what is not}

This framework:
\begin{itemize}
\item preserves AQFT locality and causal propagation;
\item explains entanglement as non-factorization induced by coherent admissibility constraints;
\item explains measurement-induced partial disentanglement as disturbance + stabilization;
\item accounts for Bell and temporal Bell violations without nonlocal influence.
\end{itemize}

It does not claim a unique global spacetime history or a global S-matrix on generic curved backgrounds.
Where asymptotic flatness fails, the correct description is via local algebras and states.

% ============================================================
\section{Conclusion}

AQFT already reconciles entanglement with locality by separating algebraic locality (commutators) from state factorization.
MTT adds a physically motivated state-selection principle: admissible coherent-sector states.
Within this restriction, entanglement is interpreted as a coherence constraint in configuration space, while microcausality remains exact.
Measurement reduces entanglement via localized disturbance and stabilization into admissible basins.
Entanglement spreads through local interactions via common ancestors and finite-speed propagation.
This supplies a coherent mathematical and physical account unifying AQFT locality with the MTT measurement and Bell analyses.

\appendix
\section{A Ten--Dimensional Formulation in Fixed--Point Language}\label{app:10D-FP}

\subsection{10D local algebras and locality}

Let $M^{10}$ denote the full Modal Triplet manifold equipped with the bundle/folding structure
introduced in the Fixed Points series. For any causally admissible region
$\mathcal{U}\subset M^{10}$, define a local observable algebra
$\mathcal{A}_{10}(\mathcal{U})$ generated by fields supported in $\mathcal{U}$.

Locality in ten dimensions is imposed in the standard algebraic sense:
for spacelike--separated regions $\mathcal{U}_1,\mathcal{U}_2\subset M^{10}$,
\begin{equation}\label{eq:10D-locality}
[\mathcal{A}_{10}(\mathcal{U}_1),\mathcal{A}_{10}(\mathcal{U}_2)]=0.
\end{equation}
This axiom is purely kinematic and independent of any coherence or stability assumptions.

\subsection{The coherent sector as an admissible invariant set}

Let $\Pi_{\mathrm{coh}}$ denote the coherent projector defined in FP~I--II, acting on the
10D configuration space $\mathcal{H}_{10}$ and selecting the joint low--lying spectral
sector compatible with the modal bundle structure.

Define the coherent configuration space
\[
\mathcal{H}_{\mathrm{coh}}:=\operatorname{Ran}(\Pi_{\mathrm{coh}}).
\]

As in FP~III--V, we restrict attention to an admissible invariant set
$D\subset\mathcal{H}_{\mathrm{coh}}$ satisfying:
\begin{enumerate}[label=\textup{(\roman*)}]
\item bounded geometry on time slabs;
\item persistence of a uniform spectral gap separating coherent from noncoherent modes;
\item boundedness of $\Pi_{\mathrm{coh}}$ on the relevant Sobolev scales;
\item disturbance--damping balance (FCC), ensuring local contractivity of the projected
      evolution.
\end{enumerate}

Physical configurations are elements of $D$. This restriction is dynamical and local; it
does not modify the algebraic locality axiom \eqref{eq:10D-locality}.

\subsection{States and non-factorization in FP language}

A ten--dimensional state is a positive normalized linear functional
\[
\omega_{10}:\mathcal{A}_{10}\to\mathbb{C}.
\]
We say $\omega_{10}$ is \emph{admissible} if its support lies in $D$.

\begin{proposition}[Global constraints without nonlocality]\label{prop:10D-nonfact}
Let $\mathcal{U}_1,\mathcal{U}_2\subset M^{10}$ be spacelike separated.
For $\omega_{10}$ admissible, expectation values need not factorize:
\[
\omega_{10}(AB)\neq \omega_{10}(A)\,\omega_{10}(B),
\quad A\in\mathcal{A}_{10}(\mathcal{U}_1),\;B\in\mathcal{A}_{10}(\mathcal{U}_2),
\]
even though \eqref{eq:10D-locality} holds.
\end{proposition}

\begin{remark}
In FP terms, this non-factorization reflects the fact that the admissible invariant set
$D$ is not a product set under decomposition into subregions. Admissibility imposes
global constraints in configuration space while preserving local spacetime causality.
\end{remark}

\subsection{Measurement as disturbance and contraction}

Measurement in the FP framework is modeled as a localized disturbance followed by
stabilization under the projected dynamics. Schematically,
\begin{equation}\label{eq:FP-measure}
\Psi \;\mapsto\; T_\tau\bigl(\Psi+\delta\Psi\bigr),
\qquad
T_\tau:=\Pi_{\mathrm{coh}}\circ\Phi_\tau,
\end{equation}
where $\delta\Psi$ is a bounded local perturbation and $\Phi_\tau$ is the unprojected
time--$\tau$ evolution.

By FCC, $T_\tau$ is contractive on $D$. Hence disturbances are damped and the configuration
returns to an admissible basin.

\begin{proposition}[Partial disentanglement]\label{prop:10D-partial}
Let $\Psi\in D$ encode correlations between degrees of freedom supported in disjoint
regions. A local disturbance acting in one region followed by \eqref{eq:FP-measure}
typically reduces cross--region correlations by contracting the configuration into a
smaller admissible basin.
\end{proposition}

This is the FP--language counterpart of partial disentanglement in the 4D CP--map
description.

\subsection{Entanglement propagation and common ancestors}

Consider three regions $\mathcal{U}_A,\mathcal{U}_B,\mathcal{U}_C\subset M^{10}$.
Suppose $\Psi$ evolves such that:
\begin{itemize}
\item $\mathcal{U}_A$ overlaps with $\mathcal{U}_C$ at some time,
\item $\mathcal{U}_B$ overlaps with $\mathcal{U}_C$ at a later time,
\item $\mathcal{U}_A$ and $\mathcal{U}_B$ never overlap directly.
\end{itemize}

\begin{proposition}[Common--ancestor correlations]\label{prop:ancestor}
After tracing out degrees of freedom in $\mathcal{U}_C$, the reduced configuration on
$\mathcal{U}_A\cup\mathcal{U}_B$ may exhibit non-factorizing correlations, even though
no direct interaction occurred between $\mathcal{U}_A$ and $\mathcal{U}_B$.
\end{proposition}

\begin{remark}
In FP terms, $\mathcal{U}_C$ acts as a \emph{common ancestor} by imposing a joint
admissibility constraint on the configuration in $D$. Correlations propagate through
local overlap and projected evolution, not by superluminal influence.
\end{remark}

\subsection{Projection to four dimensions}

The effective four--dimensional AQFT description is obtained by projecting admissible
10D states to observables supported on the physical spacetime $Y^4\subset M^{10}$.
Denote this pushforward by
\[
\omega_{4}(A):=\omega_{10}(\tilde A),
\]
where $\tilde A$ is the natural lift of $A\in\mathcal{A}(Y^4)$.

\begin{proposition}[4D entanglement as a projection artifact]\label{prop:projection}
Non-factorizing correlations in $\omega_{4}$ arise as the projection of admissible
10D coherence constraints. Local commutativity in 4D is preserved because it descends
from \eqref{eq:10D-locality}.
\end{proposition}

\subsection{Alignment with FP~I--VI}

This appendix shows that the entanglement framework used in the main text is the
direct AQFT translation of FP~I--VI principles:
\begin{itemize}
\item admissible states correspond to invariant sets $D$;
\item Bell and temporal Bell violations reflect non-product structure of $D$;
\item measurement corresponds to disturbance plus contraction under $T_\tau$;
\item entanglement propagation is causal transport of constraints through local overlap.
\end{itemize}

Thus the 4D entanglement picture is fully aligned with the ten--dimensional fixed--point
framework and introduces no additional assumptions beyond those already present in the
FP series.

% ============================================================
% References (manual, arXiv-safe)
% ============================================================
\begin{thebibliography}{99}

\bibitem{Haag1996}
R.~Haag,
\newblock {\em Local Quantum Physics: Fields, Particles, Algebras},
\newblock Springer, 2nd ed., 1996.

\bibitem{BFV2003}
R.~Brunetti, K.~Fredenhagen, and R.~Verch,
\newblock The generally covariant locality principle: A new paradigm for local quantum field theory,
\newblock {\em Commun.\ Math.\ Phys.}, 237:31--68, 2003.

\bibitem{HollandsWald2001}
S.~Hollands and R.~M.~Wald,
\newblock Local Wick polynomials and time ordered products of quantum fields in curved spacetime,
\newblock {\em Commun.\ Math.\ Phys.}, 223:289--326, 2001.

\bibitem{Bell1964}
J.~S.~Bell,
\newblock On the Einstein Podolsky Rosen paradox,
\newblock {\em Physics}, 1:195--200, 1964.

\bibitem{CHSH1969}
J.~F.~Clauser, M.~A.~Horne, A.~Shimony, and R.~A.~Holt,
\newblock Proposed experiment to test local hidden-variable theories,
\newblock {\em Phys.\ Rev.\ Lett.}, 23:880--884, 1969.

\bibitem{Tsirelson1980}
B.~S.~Tsirelson,
\newblock Quantum generalizations of Bell’s inequality,
\newblock {\em Lett.\ Math.\ Phys.}, 4:93--100, 1980.

\bibitem{LeggettGarg1985}
A.~J.~Leggett and A.~Garg,
\newblock Quantum mechanics versus macroscopic realism: Is the flux there when nobody looks?,
\newblock {\em Phys.\ Rev.\ Lett.}, 54:857--860, 1985.

\bibitem{LiebRobinson1972}
E.~H.~Lieb and D.~W.~Robinson,
\newblock The finite group velocity of quantum spin systems,
\newblock {\em Commun.\ Math.\ Phys.}, 28:251--257, 1972.

% -------------------------
% MTT / Fixed Points corpus
% -------------------------

\bibitem{MTTFoundation}
P.~Nero,
\newblock Modal Triplet Theory: Foundation --- A rigorous fixed-point framework for unified 4D physics,
\newblock Preprint, 2025.

\bibitem{FP3}
P.~Nero,
\newblock Fixed Points III: Disturbance--Damping Balance and Stability,
\newblock Preprint, 2025.

\bibitem{FP5}
P.~Nero,
\newblock Fixed Points V: Curvature Coupling and Admissibility Barriers,
\newblock Preprint, 2025.

\bibitem{FP6}
P.~Nero,
\newblock Fixed Points VI: Formal Synthesis and Physical Interpretations,
\newblock Preprint, 2025.

\bibitem{NeroMTTtoQFT}
P.~Nero,
\newblock Modal Triplet Theory: From MTT to Quantum Field Theory on Curved Spacetime,
\newblock Preprint, 2025.

\bibitem{NeroMeasurement}
P.~Nero,
\newblock Measurement as Disturbance and Stabilization in Modal Triplet Theory,
\newblock Preprint, 2025.

\bibitem{NeroTemporalBell}
P.~Nero,
\newblock Temporal Bell Inequalities and Global Consistency,
\newblock Preprint, 2025.

\bibitem{NeroBellBeables}
P.~Nero,
\newblock Modal Fixed Points, Bell's Beables, and the Limits of Factorization,
\newblock Preprint, 2025.

\end{thebibliography}

\end{document}
